\documentclass[11pt]{article}
\usepackage[margin=1in]{geometry}
\usepackage{amsmath,amssymb,amsthm}
\usepackage[T1]{fontenc}
\usepackage[utf8]{inputenc}
\usepackage{hyperref}
\usepackage{enumitem}
\usepackage{verbatim}
\usepackage{array}

\title{Plain-to-MeTTa and Plain-to-Rholang: A SpecAtom IR Proposal}
\author{ZeroBot / ProtoCosmoBot working note for Ben Goertzel}
\date{2026-06-29}

\newtheorem{principle}{Principle}
\newtheorem{proposal}{Proposal}
\newtheorem{risk}{Risk}
\newtheorem{question}{Open Question}

\begin{document}
\maketitle

\begin{abstract}
This note reviews a Telegram discussion among Ben Goertzel, ProtoCosmoBot, and
ProtoMegaTron about building a toolchain from the \texttt{***plain} software
specification language to MeTTa, initially PeTTa, and later Rholang and
MeTTa-IL. The central recommendation is to avoid a direct natural-language to
code generator. Instead, parse the structured parts of Plain into a small,
typed, Atomspace-like logical intermediate representation, called here
\emph{SpecAtom IR}. This IR should be close enough to MeTTa that MeTTa/PeTTa
lowering is natural, while preserving enough type and proof structure to
support later OSLF-style Curry-Howard correspondences and PLN-based spec
analysis. The MVP should use only shallow natural-language processing but
serious type discipline.
\end{abstract}

\section{Context and motivation}

Ben pointed to the Plain documentation and GitHub organization and asked about
the viability of a tool that maps Plain software specifications into MeTTa,
using PeTTa as the first concrete MeTTa dialect, and then into Rholang. A later
extension to MeTTa-IL was also proposed.

Plain, as described in its documentation, is a textual specification language
for spec-driven development. Its structure includes sections such as:

\begin{itemize}[nosep]
\item definitions;
\item implementation requirements;
\item test requirements;
\item functional specifications;
\item acceptance tests nested under functional specifications.
\end{itemize}

The key observation is that Plain is not just free-form prose. It already has a
sectional structure that separates concepts, implementation choices, behavior,
and tests. That structure is sufficient for a useful first compiler even if the
English content of each bullet is only minimally interpreted.

The discussion converged on the view that the project is viable, but should be
designed as a staged semantic compiler rather than a direct prose-to-code
translator:

\begin{center}
\texttt{Plain -> PlainAST -> SpecAtom IR -> MeTTa/PeTTa -> Rholang / MeTTa-IL}
\end{center}

\section{Why not compile Plain directly to MeTTa or Rholang?}

A direct compiler from Plain text to MeTTa or Rholang would be tempting, but it
would prematurely mix several distinct semantic roles:

\begin{itemize}[nosep]
\item concepts and types;
\item definitional assertions;
\item behavioral obligations;
\item implementation constraints;
\item test oracles;
\item constructive witnesses or program fragments;
\item process and concurrency constraints.
\end{itemize}

If these roles are collapsed, the resulting tool becomes mostly a code generator
with opaque prompt engineering. The more interesting possibility is to make the
specification inspectable, queryable, transformable, and analyzable before code
is emitted.

This suggests a logical IR. The IR can preserve the role of each Plain clause,
while still being easy to lower to executable or semi-executable targets.

\section{SpecAtom IR: an Atomspace-like logical mirror}

Ben suggested that the IR may best be an Atomspace-based logical representation.
There are two major reasons.

\subsection{A logical mirror of MeTTa}

If the IR is designed as a logical mirror of MeTTa, then the mapping to MeTTa
code is not an ad hoc backend. It becomes a lowering from typed logical objects
to executable or rewrite-oriented MeTTa forms. This is especially appealing if
the type relationships are arranged to support Curry-Howard-style and
Meredith/Stay OSLF-style correspondences.

The guiding idea is that specifications, types, proofs, programs, tests, and
processes should not be unrelated syntactic classes. They should be related by
explicit correspondences. For example:

\begin{center}
\begin{tabular}{>{\raggedright\arraybackslash}p{0.28\linewidth} >{\raggedright\arraybackslash}p{0.58\linewidth}}
Spec role & Logical/programmatic role \\
\hline
Concept & Type or sort \\
Functional requirement & Proposition or obligation \\
Implementation fragment & Constructive witness or program term \\
Acceptance test & Refutation attempt, observation predicate, or validation procedure \\
Process requirement & Behavioral protocol or transition-system constraint \\
Refinement relation & Proof obligation or decomposition relation \\
\end{tabular}
\end{center}

This table is not yet a formal OSLF mapping, but it indicates the kind of type
separation needed from the start.

\subsection{Future PLN-based spec analysis}

An Atomspace-like IR also opens the door to PLN-based analysis. Requirements are
often incomplete, uncertain, redundant, or inconsistent. If they are represented
as typed logical objects rather than opaque text, then later tools can ask:

\begin{itemize}[nosep]
\item Are two requirements contradictory?
\item Is a concept underspecified?
\item Does an acceptance test actually test its parent functional requirement?
\item Is an implementation requirement inconsistent with a process requirement?
\item Are there missing tests for high-risk obligations?
\item Which Plain clauses support or weaken a proposed MeTTa or Rholang artifact?
\end{itemize}

For an MVP, such PLN analysis need not be implemented. The important point is to
choose an IR that does not block it.

\section{A small typed schema for the MVP}

The main danger of an Atomspace-like IR is overgeneralization. A full Atomspace
ontology could consume the project before a working prototype exists. The MVP
should therefore define a deliberately small schema, perhaps called
\texttt{SpecAtom}. The schema should be narrow enough for deterministic parsing
from Plain sections but typed enough to support later logical correspondences.

A possible initial vocabulary is:

\begin{verbatim}
(PlainFile file-id path)
(Section section-id file-id SectionKind)
(PlainItem item-id section-id ordinal raw-text)

(Concept concept-id name)
(Defines item-id concept-id raw-text)
(Attribute concept-id attribute-id Requiredness ValueKind raw-text)

(ImplementationReq req-id item-id raw-text)
(TestReq req-id item-id raw-text)
(FunctionalReq req-id item-id raw-text)
(AcceptanceTest test-id parent-functional-req-id raw-text)

(Refines child-req-id parent-req-id)
(DependsOn req-id other-req-id)
(Constrains constraint-id target-id raw-text)

(Obligation req-id proposition-id)
(Witness witness-id obligation-id target-language artifact-ref)
(TestOracle test-id predicate-id)
(ProcessConstraint constraint-id protocol-id)
\end{verbatim}

This is intentionally modest. The raw English remains present. The compiler does
not pretend to understand all English. However, each artifact receives a logical
role and a place in the type structure.

\begin{principle}[Role preservation]
Every generated artifact should know whether it is a type or concept, a
proposition/spec obligation, a constructive witness/program fragment, a
test/refutation condition, or a process/interaction constraint.
\end{principle}

\begin{principle}[Raw-text provenance]
Every logical object generated from Plain should retain a pointer to the source
file, section, bullet, and raw text span that produced it.
\end{principle}

\begin{principle}[Shallow NLP first]
The MVP should avoid full English-to-logic parsing. It should parse Plain's
explicit structure and concept references, then preserve the remaining prose as
raw text attached to typed logical objects.
\end{principle}

\section{Mapping Plain to SpecAtom IR}

The Plain-to-IR compiler can be mostly deterministic at first.

\subsection{Definitions}

Plain definitions introduce concepts and possibly attributes. For example:

\begin{verbatim}
***definitions***
- :Task: describes an activity that needs to be done by :User:. :Task:
  has the following attributes:
  - Name - a short description of :Task:. This is a required attribute.
  - Due Date - optional date by which :User: is supposed to complete :Task:.
\end{verbatim}

A shallow compiler can emit:

\begin{verbatim}
(Concept :Task "Task")
(Concept :User "User")
(Defines item-17 :Task "describes an activity ...")
(Attribute :Task :Name Required Text "a short description ...")
(Attribute :Task :DueDate Optional Date "optional date ...")
\end{verbatim}

The exact attribute typing can begin heuristic-based and become more precise
later.

\subsection{Implementation requirements}

Plain implementation requirements describe how the software should be built.
They should not be confused with behavioral requirements. For example:

\begin{verbatim}
- :Implementation: should be in Python.
\end{verbatim}

This may become:

\begin{verbatim}
(ImplementationReq impl-1 item-21 ":Implementation: should be in Python.")
(Constrains impl-constraint-1 :Implementation "target-language=Python")
\end{verbatim}

The first MVP can store most of this as typed raw text. Later versions can
recognize standard implementation templates.

\subsection{Functional specifications}

Functional specs describe behavior. A bullet such as:

\begin{verbatim}
- :User: should be able to add :Task:. Only valid :Task: items can be added.
\end{verbatim}

can become:

\begin{verbatim}
(FunctionalReq req-33 item-33 raw-text-33)
(Obligation req-33 prop-33)
(References req-33 :User)
(References req-33 :Task)
\end{verbatim}

A later semantic pass may infer action structure:

\begin{verbatim}
(ActionSchema action-33 :User add :Task)
(Precondition action-33 valid-task)
\end{verbatim}

But the MVP does not need to rely on that inference.

\subsection{Acceptance tests}

Acceptance tests should be attached to their parent functional requirement and
represented as test oracles or observation predicates:

\begin{verbatim}
(AcceptanceTest test-34 req-33 raw-test-text)
(TestOracle test-34 predicate-34)
(Tests test-34 req-33)
\end{verbatim}

This is important because it lets later analysis ask whether a test actually
covers its parent obligation.

\section{Mapping SpecAtom IR to MeTTa/PeTTa}

The first backend should be MeTTa/PeTTa because the IR is naturally atom-shaped.
There are two levels of output.

\subsection{Reified specification atoms}

The simplest output is a faithful MeTTa representation of the IR:

\begin{verbatim}
(: Task Concept)
(Defines item-17 Task "describes an activity ...")
(Attribute Task Name Required Text)
(FunctionalReq req-33 item-33)
(References req-33 User)
(References req-33 Task)
(AcceptanceTest test-34 req-33)
\end{verbatim}

This is already useful. It lets MeTTa queries inspect the specification, trace
requirements, and drive later transformations.

\subsection{Executable or semi-executable PeTTa skeletons}

A second pass can lower selected functional requirements into executable
skeletons. For example, if a functional requirement is recognized as an action,
the compiler may generate PeTTa rules or stubs such as:

\begin{verbatim}
(: add-task (-> User Task TaskList TaskList))
(= (add-task $u $task $list)
   (if (valid-task $task)
       (append-task $task $list)
       (error InvalidTask)))
\end{verbatim}

For the MVP, these skeletons can be conservative. If the compiler is unsure, it
should emit a typed obligation plus a TODO witness rather than hallucinating an
implementation.

\section{Mapping SpecAtom IR to Rholang}

Rholang is a less direct target than MeTTa because it is process-oriented. The
IR should only lower to Rholang when the Plain specification includes enough
process-like structure:

\begin{itemize}[nosep]
\item services or agents;
\item messages or events;
\item channels or resources;
\item state transitions;
\item concurrency or consensus constraints;
\item protocols with observable traces.
\end{itemize}

The Rholang backend should therefore be a later stage, not the first backend.
It can consume IR objects such as:

\begin{verbatim}
(Process :TaskService)
(Channel :AddTaskRequest)
(Channel :TaskAddedEvent)
(Transition add-task-transition pre-state message post-state)
(ProcessConstraint c-12 protocol-7)
\end{verbatim}

and emit Rholang process templates. Requirements that are purely functional or
local should remain in MeTTa/PeTTa or be represented as comments/obligations in
Rholang until a process interpretation is supplied.

\section{OSLF-style type relationships}

Ben emphasized that even without full English-to-logic parsing, the type
relationships should be correct so that OSLF-style correspondences work later.
This suggests the following initial type hierarchy.

\begin{verbatim}
SpecObject
  SourceObject
    PlainFile
    Section
    PlainItem
  LogicalObject
    Concept
    Proposition
    Obligation
    Constraint
    Refinement
  ConstructiveObject
    Witness
    ProgramFragment
    MeTTaTerm
    RholangProcess
  ValidationObject
    AcceptanceTest
    TestOracle
    ObservationPredicate
  ProcessObject
    Channel
    Message
    Protocol
    Transition
\end{verbatim}

The key correspondence pattern is:

\begin{center}
\begin{tabular}{>{\raggedright\arraybackslash}p{0.3\linewidth} >{\raggedright\arraybackslash}p{0.55\linewidth}}
Type-theoretic role & SpecAtom role \\
\hline
Type & Concept or interface sort \\
Term & Program fragment, action, or process term \\
Proposition & Functional requirement or constraint \\
Proof/witness & Implementation artifact satisfying an obligation \\
Counterexample & Failed acceptance test, trace violation, or model check failure \\
Refinement & Decomposition of an obligation into smaller obligations \\
\end{tabular}
\end{center}

This is enough to guide the MVP without forcing a complete formal semantics.

\section{PLN-based analysis enabled later}

Once Plain specs are represented as SpecAtom IR, PLN-style reasoning can operate
on the specification graph. Examples include:

\begin{itemize}[nosep]
\item estimating whether a functional requirement is supported by acceptance tests;
\item detecting probable contradictions among requirements;
\item identifying concepts used without definitions;
\item suggesting refinements for overly complex Plain functional specs;
\item propagating confidence from tests to obligations and witnesses;
\item comparing two implementations as alternative witnesses for the same
obligation;
\item asking whether a Rholang process preserves the obligations of the MeTTa
specification.
\end{itemize}

The MVP does not need PLN. It needs to avoid destroying the information that PLN
would require later.

\section{Proposed MVP}

A realistic MVP could be:

\begin{enumerate}
\item Parse a subset of Plain files: section headers, bullet hierarchy, concept
notation, and acceptance-test nesting.
\item Emit a JSON or S-expression PlainAST for debugging.
\item Emit SpecAtom IR in MeTTa syntax.
\item Provide simple MeTTa queries: list concepts, list obligations, list tests
for an obligation, list undefined concept references.
\item Generate a conservative PeTTa skeleton for one small example, such as a
hello-world app or task-manager fragment.
\item Preserve all raw text and source locations.
\item Reject or mark as underspecified anything beyond the shallow subset.
\end{enumerate}

The MVP should not attempt full English parsing, automatic correct code
synthesis, or general Plain-to-Rholang generation.

\section{Development plan}

\subsection{Phase 0: corpus and parser}

Collect a few Plain examples from the public examples repository and hand-write
one or two examples aimed at MeTTa/PeTTa. Implement a parser for Plain's section
structure and bullet hierarchy.

\subsection{Phase 1: SpecAtom schema}

Define the small SpecAtom vocabulary. Write schema validation rules so that each
artifact has exactly one primary role and a source pointer.

\subsection{Phase 2: MeTTa/PeTTa backend}

Emit faithful reified specification atoms. Add simple query examples and a small
PeTTa skeleton generator for recognized action patterns.

\subsection{Phase 3: type correspondence checks}

Add checks inspired by Curry-Howard/OSLF:

\begin{itemize}[nosep]
\item every witness targets an obligation;
\item every test targets an obligation;
\item every obligation references defined concepts or explicitly external
concepts;
\item refinements preserve parent-child role consistency;
\item process constraints are not lowered to Rholang unless process roles are
present.
\end{itemize}

\subsection{Phase 4: Rholang backend}

Add a restricted process sublanguage and emit Rholang templates only from that
sublanguage.

\subsection{Phase 5: PLN analysis}

Layer PLN-style uncertainty and support relations over the SpecAtom graph. Use
this for contradiction detection, coverage estimates, and refinement
recommendations.

\section{Risks and mitigations}

\begin{risk}[Overbuilt ontology]
A full Atomspace ontology could delay a working compiler indefinitely.
\end{risk}

\noindent Mitigation: define a tiny SpecAtom subset and require end-to-end
Plain-to-MeTTa output in the first prototype.

\begin{risk}[False semantic precision]
The compiler may appear to understand English that it only stores as raw text.
\end{risk}

\noindent Mitigation: mark shallowly parsed clauses explicitly. Use confidence,
source pointers, and TODO witnesses rather than pretending full semantic
interpretation.

\begin{risk}[Rholang mismatch]
Many Plain specs may not contain enough process structure for meaningful
Rholang output.
\end{risk}

\noindent Mitigation: require explicit process annotations before Rholang
lowering.

\begin{risk}[Loss of OSLF compatibility]
If early shortcuts collapse obligations, witnesses, tests, and processes, later
OSLF-style mappings become awkward.
\end{risk}

\noindent Mitigation: enforce role preservation from the start, even if the
English interpretation is shallow.

\section{Conclusion}

The proposed Plain-to-MeTTa and Plain-to-Rholang project is viable and
potentially valuable, but its core should be a typed Atomspace-like IR rather
than direct code generation. The simplest useful prototype should parse Plain's
explicit structure, preserve raw text and provenance, generate typed SpecAtom
objects, and lower those objects to MeTTa/PeTTa. This is sufficient to make Plain
specs inspectable and queryable while preserving a path toward OSLF-style type
correspondences and PLN-based analysis.

The guiding maxim is: keep the first compiler shallow in natural-language
understanding but deep in type discipline.

\end{document}
